% ZTL Blueprint: formalization map of the core. All \leanok items are proven in
% the repository's Lean corpus (lean/) with an empty axiom list. Updated
% 2026-10-30 against the core as of preprint 2.1.6 (judge vocabulary of
% 2026-10-29). Every name in a lean tag is checked against the built corpus.

\chapter*{About this blueprint}

This blueprint is a dependency graph of the \emph{core} of ZTL: the values and
the one rule that generates every connective, the sixteen lifted tables, the
agreement with classical logic on verified inputs, the twenty-six laws and the
rules of inference, the two registers and the judge's dispositions, the
relation to Bochvar's external logic, and the receipt. Each node names the Lean
declarations that prove it; it does not try to cover the whole corpus (about
2800 theorems, most of them about numbers and applications). The human reading
of the same material is on the site:
\href{https://ztljudge.com/ztl/rule/}{the rule} and
\href{https://ztljudge.com/ztl/theorems/}{the theorems}; the preprint is
\href{https://doi.org/10.5281/zenodo.23776523}{doi:10.5281/zenodo.23776523}.

\chapter{The values and the rule}

\begin{definition}[Values and the mark]\label{def:values}
  \lean{V}\leanok
  Truth values are $T$ and $F$; the third symbol $Z$ is a mark for a dirty
  (unverified) input, not a truth value.
\end{definition}

\begin{definition}[Generating principle: the lift]\label{def:lift}
  \uses{def:values}\lean{V.subs, V.lift1, V.lift2}\leanok
  $Z$ has two classical readings, $T$ has one, $F$ has one. For a classical
  connective $f$, its ZTL lift returns $T$ iff $f$ is $T$ under \emph{every}
  classical reading of its arguments (each occurrence of $Z$ read on its own);
  otherwise $F$. Truth is never granted on credit.
\end{definition}

\begin{definition}[Connectives]\label{def:ops}
  \uses{def:lift}\lean{V.znot, V.zand, V.zor, V.zimp, V.zxor, V.zxnor}\leanok
  All six connectives are lifts of their classical cores; no table is
  postulated.
\end{definition}

\begin{theorem}[Anchor cells]\label{thm:anchors}
  \uses{def:ops}\lean{V.ax_not_Z, V.ax_notnot_Z, V.ax_xnor_ZZ, V.ax_xnor_ZT, V.ax_xor_ZZ, V.ax_xor_ZF, V.ax_xor_ZT}\leanok
  The design anchor cells are theorems of the lift: $\neg Z=F$,
  $\neg\neg Z=T$, $Z\leftrightarrow Z=F$, $Z\leftrightarrow T=F$,
  $Z\oplus Z=Z\oplus F=Z\oplus T=F$.
\end{theorem}

\begin{theorem}[Greediness]\label{thm:greedy}
  \uses{def:ops}\lean{V.lift1_classical, V.lift2_classical, V.evalF_classical}\leanok
  Every lifted connective answers $T$ or $F$; hence every compound formula is
  classical under every valuation, over the whole formula language: $Z$ lives
  only on atoms.
\end{theorem}

\begin{theorem}[The lift is the only decisive, sound and maximal table]\label{thm:liftchar}
  \uses{def:lift}\lean{V.lift1_characterized, V.lift2_characterized}\leanok
  Call a table for a classical connective $g$ \emph{decisive} if it never
  answers $Z$, \emph{sound} if its $T$ is forced by every reading of its
  arguments, \emph{maximal} if it answers $T$ whenever every reading forces
  it. For unary and binary $g$, a table has all three properties iff it is
  the ZTL lift of $g$.
\end{theorem}

\chapter{The sixteen tables}

\begin{definition}[Every binary kernel, lifted]\label{def:sixteen}
  \uses{def:lift}\lean{mkF, lf, znand, znor, znimp}\leanok
  \texttt{mkF} enumerates the sixteen binary Boolean functions by their four
  outputs; \texttt{lf} is the lift of each. The six connectives of
  Definition~\ref{def:ops} and the lifted NAND, NOR and nonimplication are
  members of this family.
\end{definition}

\begin{theorem}[Lifts of the negated connectives]\label{thm:neglifts}
  \uses{def:sixteen, def:ops}\lean{V.nand_lift, V.nor_lift, V.nimp_lift, V.nxor_lift, V.lift_id_is_up, V.and_self_is_up}\leanok
  The lift of a negated classical connective is the argument-side De Morgan
  form: $\mathrm{lift}(\uparrow)(a,b)=\neg a\vee\neg b$,
  $\mathrm{lift}(\downarrow)(a,b)=\neg a\wedge\neg b$,
  $\mathrm{lift}(\not\to)(a,b)=a\wedge\neg b$, and the lift of
  $\neg\oplus$ is $\leftrightarrow$. The lift of the identity is the
  assertion $\uparrow x$ ($T\mapsto T$, else $F$), definable as $x\wedge x$.
\end{theorem}

\begin{theorem}[Composite $\neq$ primitive]\label{thm:composite}
  \uses{def:sixteen}\lean{negated_lift_differs, negated_lift_constants, two_nands}\leanok
  For every non-constant binary function $f$ (fourteen of the sixteen), the
  lift of $\neg f$ differs from the negated verdict $\neg\,\mathrm{lift}(f)$
  in some cell; for the two constants they agree. Example: lifted NAND gives
  $F$ at $(Z,T)$ while $\neg(Z\wedge T)=T$. So a table printed as
  $\neg(a\circ b)$ is a composition of two lifts, not the lift of the negated
  connective.
\end{theorem}

\begin{theorem}[Single-operator completeness]\label{thm:clone}
  \uses{def:sixteen, thm:neglifts}\lean{directional_complete, nondirectional_incomplete, nand_cannot_and, nand_cannot_or, nor_cannot_and, nor_cannot_or, nimp_clone_eq, imp_clone_eq}\leanok
  With the constants, the lift of a binary function generates the canonical
  clone $\{\neg,\wedge,\vee\}$ alone iff the function is one of the four
  directional ones (both implications, both nonimplications). The other
  twelve each miss negation or conjunction; in particular lifted NAND and NOR,
  complete classically, are not complete here.
\end{theorem}

\chapter{Agreement with classical logic}

\begin{definition}[Classical evaluation]\label{def:cval}
  \uses{def:values}\lean{V.emb, V.cval}\leanok
  A classical valuation is embedded into ZTL values without the mark;
  \texttt{cval} is ordinary two-valued evaluation of the same formulas.
\end{definition}

\begin{theorem}[Conservativity on verified inputs]\label{thm:conservative}
  \uses{def:cval, def:ops}\lean{V.evalF_agrees}\leanok
  On a valuation that carries no mark, ZTL evaluation and classical
  evaluation agree on every formula of the language (structural induction,
  not enumeration). In particular all twenty-six classical laws of
  Chapter~\ref{chap:laws} hold on verified data.
\end{theorem}

\begin{theorem}[No new laws]\label{thm:taut}
  \uses{thm:conservative}\lean{V.ztl_taut_is_classical}\leanok
  A formula that is $T$ under every ZTL valuation is a classical tautology.
\end{theorem}

\begin{theorem}[Strictly fewer]\label{thm:notconversely}
  \uses{thm:taut}\lean{V.not_conversely}\leanok
  The inclusion is proper: $p\to p$ is a classical tautology and is $F$ at a
  marked $p$. ZTL is conservative over classical logic, not the same logic.
\end{theorem}

\chapter{Laws and rules}\label{chap:laws}

\begin{theorem}[The twelve laws that hold on a dirty atom]\label{thm:laws12}
  \uses{def:ops}\lean{V.imp_def, V.xor_def, V.xnor_def, V.xnor_imps, V.and_comm', V.or_comm', V.and_assoc', V.or_assoc', V.and_distrib, V.or_distrib, V.no_contradiction, V.imp_trans_law}\leanok
  Definitions of $\to,\oplus,\leftrightarrow$ by $\neg,\wedge,\vee$;
  $\leftrightarrow$ as two implications; commutativity and associativity of
  $\wedge,\vee$; both distributive laws; non-contradiction
  $\neg(p\wedge\neg p)=T$; the transitivity law. Each holds at every value,
  the mark included.
\end{theorem}

\begin{theorem}[The fourteen laws that need ground]\label{thm:laws14}
  \uses{def:ops}\lean{V.double_neg_needs_ground, V.deMorgan1_needs_ground, V.deMorgan2_needs_ground, V.contraposition_needs_ground, V.xor_as_not_xnor_needs_ground, V.and_idem_needs_ground, V.or_idem_needs_ground, V.absorption_needs_ground, V.and_true_needs_ground, V.or_false_needs_ground, V.lem_needs_ground, V.imp_refl_needs_ground, V.peirce_needs_ground, V.k_axiom_needs_ground}\leanok
  Double negation, both De Morgan laws, contraposition as an identity,
  $\oplus$ as $\neg\leftrightarrow$, both idempotences, absorption, the two
  unit laws, excluded middle, $p\to p$, Peirce and the K axiom are refuted at a
  marked atom (each with a counterexample, not a silence); on verified data
  they hold (Theorem~\ref{thm:conservative}).
\end{theorem}

\begin{theorem}[Modus ponens]\label{thm:mp}
  \uses{def:ops}\lean{V.modus_ponens}\leanok
  From $p=T$ and $p\to q=T$ follows $q=T$: earned truth is transported.
\end{theorem}

\begin{theorem}[Rules of inference]\label{thm:rules}
  \uses{thm:mp}\lean{V.rule_modus_tollens, V.rule_contraposition, V.rule_disj_syllogism, V.rule_and_intro, V.rule_and_elim, V.rule_or_intro, V.rule_dn_intro, V.rule_transitivity, V.rule_K, V.rule_explosion, V.rule_resolution}\leanok
  With designated value $\{T\}$: modus tollens, contraposition, disjunctive
  syllogism, $\wedge$-introduction and elimination, $\vee$-introduction,
  $\neg\neg$-introduction, transitivity, K, explosion and resolution are sound
  rules at every value.
\end{theorem}

\begin{theorem}[Rules vs laws split]\label{thm:split}
  \uses{thm:rules, thm:laws14}\lean{V.rule_contraposition, V.rule_dn_elim_needs_ground, V.rule_taut_concl_needs_ground, V.dt_one_way}\leanok
  Contraposition survives as an inference rule while failing as an identity;
  $\neg\neg$-elimination and a tautology in the conclusion fail even as rules.
  The deduction theorem holds in one direction only: $p\vDash p$, yet
  $Z\to Z=F$.
\end{theorem}

\begin{theorem}[No gluts]\label{thm:nogluts}
  \uses{def:ops}\lean{V.no_gluts}\leanok
  $v(p)=v(\neg p)=T$ is unsatisfiable: ZTL is paracomplete, not
  paraconsistent.
\end{theorem}

\chapter{Two registers and the judge}

\begin{definition}[The cautious register: strong Kleene]\label{def:kleene}
  \uses{def:values}\lean{V.knot, V.kand, V.kor, V.k1, V.k2, V.kev, V.k_matches_core}\leanok
  The lazy register is strong Kleene: $T$ if every reading gives $1$, $F$ if
  every reading gives $0$, else $Z$. \texttt{kev} evaluates a formula in it.
\end{definition}

\begin{theorem}[Kleene register]\label{thm:kleene}
  \uses{def:kleene, thm:liar}\lean{V.kleene_and_monotone, V.eager_and_not_monotone, V.liar_kleene_home}\leanok
  The lazy register is monotone in the information order and homes the liar
  ($\neg Z=Z$); the decisive (eager) register is not monotone: grounding and
  verdicts are done in two different registers.
\end{theorem}

\begin{theorem}[The dirty reading is strong Kleene]\label{thm:dirty}
  \uses{def:kleene, def:ops}\lean{V.rd_is_kleene, V.covers_rd, V.refusal_witnessed, V.never_stricter, V.over_grant}\leanok
  Read every occurrence of a dirty atom on its own; the resulting semantics is
  strong Kleene. Every $F$ of ZTL is witnessed by such a reading, and ZTL is
  never stricter than it; the converse fails: $\neg\neg p$ is granted on a
  dirty $p$.
\end{theorem}

\begin{theorem}[Each connective is Kleene followed by the assertion]\label{thm:upkleene}
  \uses{def:kleene, thm:neglifts}\lean{V.not_is_up_kleene, V.and_is_up_kleene, V.or_is_up_kleene, V.imp_is_up_kleene, V.xor_is_up_kleene, V.xnor_is_up_kleene}\leanok
  Every ZTL connective equals the strong Kleene connective followed by the
  assertion $\uparrow$ ($Z\mapsto F$).
\end{theorem}

\begin{theorem}[ZTL is the compositional Kleene verdict]\label{thm:uniquecomp}
  \uses{thm:upkleene}\lean{V.kleene_end_verdict_not_compositional, V.ztl_unique_compositional, V.ztl_is_compositional_kleene_verdict, V.verdict_everywhere_excludes_de_morgan}\leanok
  Kleene with one verdict at the end is not compositional. Any compositional
  semantics that gives Kleene's verdict on every single connective applied to
  atoms agrees with ZTL on every formula, and no such semantics satisfies De
  Morgan.
\end{theorem}

\begin{theorem}[Two registers, three properties]\label{thm:registers}
  \uses{thm:liftchar, def:kleene}\lean{V.kleene_sound_maximal2, V.kleene_not_decisive, V.bochvar_not_sound}\leanok
  Strong Kleene is sound and maximal but not decisive ($\neg Z=Z$); Bochvar's
  external connectives are decisive but not sound ($\neg Z=T$, $Z\to F=T$);
  only the ZTL lift has all three (Theorem~\ref{thm:liftchar}).
\end{theorem}

\begin{theorem}[What Kleene cannot say]\label{thm:kleenemute}
  \uses{def:kleene, def:jops}\lean{V.kleene_cannot_jump, V.kleene_cannot_assert, V.kleene_cannot_say_established, V.kleene_cannot_say_mark, V.ztl_says_evidence}\leanok
  No strong Kleene formula says ``$x$ is established true'', ``$x$ is
  established'' or ``$x$ is a mark''; ZTL says all three with its own
  connectives: $p\wedge p$, $p\vee\neg p$, $\neg(p\leftrightarrow p)$.
\end{theorem}

\begin{theorem}[Expressive completeness]\label{thm:expcomplete}
  \uses{def:jops, thm:unary}\lean{V.Cpl.dnf_spec, V.Cpl.ztl_expressively_complete}\leanok
  Every two-valued function of a valuation that depends only on a finite list
  of atoms is the value of a ZTL formula (the J-DNF over those atoms, with
  $3^{|\mathrm{atoms}|}$ leaves).
\end{theorem}

\begin{definition}[Grades and dispositions]\label{def:grades}
  \uses{def:kleene, thm:dirty}\lean{V.refines, V.Earned, V.Refuted, V.StrictEarned, V.StrictRefuted, V.total, V.Enum, V.Speak.StableF, V.Speak.EarnedIfStable, V.Speak.RefutedIfStable, V.Speak.OnCredit, V.Speak.Open}\leanok
  The hereditary grade: $T$ under every refinement of the marks. A
  completion (cleaning) is a total refinement; \texttt{Enum} is ``$T$ in
  every completion''. The judge's dispositions (2026-10-29):
  EARNED / REFUTED = strong Kleene $T$ / $F$;
  EARNED IF STABLE / REFUTED IF STABLE = not settled by Kleene, but $T$ / $F$
  in every completion (each dirty atom read the same at all its
  occurrences);
  ON CREDIT = $T$ now, the completions disagree;
  OPEN = not $T$ now, the completions disagree.
  The judge's seventh answer, E (the claim rests on something that has no
  subject), is assigned by declaration in \texttt{ztljudge.py} and is not
  formalized in the Lean corpus.
\end{definition}

\begin{theorem}[EARNED and REFUTED are strong Kleene]\label{thm:strictkleene}
  \uses{def:grades}\lean{V.earned_of_kleene, V.strict_earned_iff_kleene, V.never_laxer, V.strict_refuted_iff_kleene}\leanok
  The strict dispositions are exactly the strong Kleene values $T$ and $F$:
  the judge answers the whole-formula question as Kleene does, no more and no
  less.
\end{theorem}

\begin{theorem}[The grant is classical enumeration]\label{thm:grant}
  \uses{def:grades, thm:strictkleene, thm:conservative}\lean{V.Speak.grant_iff_enum, V.Speak.grant_strict}\leanok
  The judge grants (EARNED or EARNED IF STABLE) exactly when the formula is
  $T$ in every completion of its dirty atoms. The grant is strict: on a dirty
  $p$ it grants neither $p$ nor $\neg p$.
\end{theorem}

\begin{theorem}[Dominance]\label{thm:dominance}
  \uses{thm:grant, thm:dirty}\lean{V.substitution_of_earned, V.strict_any, V.earned_dominates_enum, V.enum_strict, V.Speak.grant_right_where_substitution_right}\leanok
  Criterion: a grant is right iff the dirty reading earns it (strong Kleene
  $T$). Under it the hereditary grade strictly dominates the grant
  (enumeration): their right grants coincide, every wrong grant of the
  hereditary grade is a wrong grant of enumeration, and enumeration grants
  $p\to p$ on a dirty $p$, which the grade refuses. The grant in turn is right
  wherever any single classical substitution of the dirty atoms is right, and
  every substitution grants one of $p,\neg p$ that the judge refuses.
\end{theorem}

\begin{theorem}[The limits of the hereditary grade]\label{thm:canerr}
  \uses{thm:dominance}\lean{V.judge_can_err, V.Speak.lnc_earned, V.Speak.lnc_if_stable}\leanok
  The hereditary grade can be wrong under the criterion: $\neg(p\wedge\neg p)$
  is $T$ under every refinement of a dirty $p$, yet strong Kleene does not
  give it $T$; the judge reports it as EARNED IF STABLE.
\end{theorem}

\begin{theorem}[The dispositions cover every case]\label{thm:partition}
  \uses{def:grades}\lean{V.Speak.dispositions_partition}\leanok
  At every valuation each formula receives at least one of the six
  dispositions. (Coverage only: the theorem states the disjunction, not
  mutual exclusivity.)
\end{theorem}

\begin{theorem}[The judge speaks ZTL]\label{thm:speaks}
  \uses{def:grades, thm:expcomplete}\lean{V.Speak.allR_iff, V.Speak.allC_iff, V.Speak.judge_speaks_ztl}\leanok
  For a fixed formula $\varphi$, each of the six dispositions is decided by
  strong Kleene and a finite search over the completions of its dirty atoms,
  and is the $T$-value of a ZTL formula over the atoms of $\varphi$. The
  judge is not code over ZTL but ZTL about itself (expressible, not cheap: the
  J-DNF has $3^{|\mathrm{atoms}|}$ leaves).
\end{theorem}

\chapter{Bochvar's external logic}

\begin{definition}[$B^e$ and the translation $\zeta$]\label{def:bochvar}
  \uses{def:values}\lean{V.Boch.BFm, V.Boch.beval, V.Boch.zeta}\leanok
  $B^e$: weak Kleene $\neg,\wedge,\vee$ ($Z$ contaminates), the external
  operators $J_0,J_1,J_2$ and the constants, designated value $T$. $\zeta$
  sends $\neg$ to $J_0$ and each binary connective to its classical form
  over $J_0,J_2$ of the arguments.
\end{definition}

\begin{theorem}[$\zeta$ is exact]\label{thm:zeta}
  \uses{def:bochvar, def:ops}\lean{V.Boch.zeta_exact, V.Boch.zeta_consequence}\leanok
  $\zeta\varphi$ has, at every valuation, the value $\varphi$ has in ZTL;
  hence $\Gamma\vDash\varphi$ in ZTL iff $\zeta\Gamma\vdash\zeta\varphi$ in
  $B^e$.
\end{theorem}

\begin{theorem}[ZTL is the two-valued fragment of $B^e$]\label{thm:externalztl}
  \uses{thm:zeta, thm:expcomplete}\lean{V.Boch.external_is_ztl}\leanok
  Every $B^e$ formula that never takes the middle value (every external
  formula among them) has, at every valuation, the value of some ZTL formula.
  By the functions they compute, ZTL is the two-valued fragment of $B^e$ in
  another basis, up to bare atoms (which carry $Z$ in both).
\end{theorem}

\begin{theorem}[Bochvar is the default ``dirty := false'']\label{thm:bochdefault}
  \uses{def:grades, thm:registers}\lean{V.bochvar_is_default_false, V.earned_dominates_bochvar, V.bochvar_strict}\leanok
  Bochvar's external connectives compute the classical value with every dirty
  atom set to $F$. What the hereditary grade grants, Bochvar grants; on a
  dirty $p$, Bochvar grants $\neg p$, which the grade and the judge refuse.
\end{theorem}

\chapter{The receipt}

\begin{definition}[The label]\label{def:label}
  \uses{def:kleene}\lean{V.labF}\leanok
  With every value the judge returns a label: the dirty atoms still holding the
  answer up.
\end{definition}

\begin{theorem}[The receipt is complete]\label{thm:receipt}
  \uses{def:label, thm:greedy}\lean{V.receipt_complete, V.receipt_complete_greedy}\leanok
  For every formula, valuation and dirty atom: if the label leaves the atom
  out, setting it to anything changes neither the lazy value nor the greedy
  verdict. An atom that could change the answer is always on the receipt. The
  label is a sound over-approximation, not exact in general.
\end{theorem}

\begin{theorem}[Exact on linear claims]\label{thm:labelexact}
  \uses{thm:receipt}\lean{V.linear_no_loss, V.drivable, V.label_exact_linear, V.clash_names_an_idle_atom}\leanok
  Where each dirty atom occurs at most once, the judge loses no truth and the
  label names only atoms that can move the verdict jointly with the others.
  Without linearity this fails: $p\wedge\neg p$ names $p$, which no reading
  moves.
\end{theorem}

\chapter{Tableau calculus, certified}

\begin{definition}[Signs]\label{def:signs}
  \uses{def:values}\lean{V.SignT, V.SignF, V.SignP, V.SignN}\leanok
  Strict signs $T,F$ and weak signs $P=\{T,Z\}$, $N=\{F,Z\}$; weak signs
  occur only in refutation polarity.
\end{definition}

\begin{definition}[Engine]\label{def:closes}
  \uses{def:signs, thm:greedy}\lean{V.closes}\leanok
  Fuel-indexed structural tableau engine over the generating basis
  $\{\neg,\wedge,\vee\}$; heavy connectives reduce by the living
  identities.
\end{definition}

\begin{theorem}[Soundness and completeness]\label{thm:closesiff}
  \uses{def:closes}\lean{V.closes_iff}\leanok
  The tableau closes iff the signed nodes are unsatisfiable --- for all
  formulas, by induction on weighted size. Zero axioms.
\end{theorem}

\begin{theorem}[Entailment certificate]\label{thm:tproves}
  \uses{thm:closesiff}\lean{V.tproves_iff}\leanok
  $\Gamma\vdash\varphi$ by the engine iff every valuation making all
  premises $T$ makes $\varphi$ $T$.
\end{theorem}

\begin{theorem}[Native engine agrees]\label{thm:agree}
  \uses{thm:tproves}\lean{V.closesN_iff, V.engines_agree}\leanok
  The engine with native signed rules for $\to,\oplus,\leftrightarrow$
  is certified by the same induction and returns identical verdicts.
\end{theorem}

\chapter{Registers and paradoxes}

\begin{theorem}[Liar is homeless]\label{thm:liar}
  \uses{def:ops}\lean{V.liar_homeless}\leanok
  Negation has no fixed point: no value seats the liar; quarantine is
  irremovable.
\end{theorem}

\begin{theorem}[Carousel and revenge]\label{thm:carousel}
  \uses{thm:kleene}\lean{V.carousel_no_model, V.revenge_bullet}\leanok
  The two-sentence cycle has no eager model over all nine value pairs;
  the revenge sentence's content evaluates to $T$ while the sentence
  stays quarantined --- the price is paid explicitly.
\end{theorem}

\begin{theorem}[Dynamics]\label{thm:dynamics}
  \uses{thm:liar}\lean{V.liar_period2, V.carousel_period4, V.curry_homeless, V.yablo3_unique}\leanok
  Liar oscillates with period 2, the carousel with period 4; Curry is
  homeless without using negation; the finite Yablo truncation has a
  unique grounded model.
\end{theorem}

\chapter{Sets with unverified elements}

\begin{definition}[Marked elements]\label{def:el}
  \uses{def:values}\lean{V.El, V.eqAtom}\leanok
  Elements are verified values or marks; the equality atom yields $Z$
  whenever a mark is involved --- even a mark with itself.
\end{definition}

\begin{theorem}[Marks earn no membership]\label{thm:memz}
  \uses{def:el}\lean{V.memZ}\leanok
  $\mathrm{mem}(z, S)=F$ for every mark $z$ and every list $S$:
  membership is an existential fold with strict witnesses.
\end{theorem}

\begin{theorem}[Identity laws fall on marked sets]\label{thm:marked}
  \uses{thm:memz}\lean{V.sub_marked_false, V.seteq_self_marked}\leanok
  A marked set is not provably a subset of itself; reflexivity of set
  equality fails --- inherited from the tables, not postulated.
\end{theorem}

\chapter{The algebraic passport}

\begin{definition}[J-operators]\label{def:jops}
  \uses{def:ops}\lean{V.jT, V.jF}\leanok
  $J_T(p) = p\wedge p$ (fallen idempotence is the truth detector),
  $J_F(p) = \neg p\wedge\neg p$, $J_Z = \mathrm{isZ}$: three exact,
  disjoint value indicators grown from ZTL's own connectives.
\end{definition}

\begin{theorem}[Expressive completeness, unary]\label{thm:unary}
  \uses{def:jops}\lean{V.unary_complete}\leanok
  Every unary external function $V\to\{T,F\}$ is a J-DNF: one theorem
  over all eight target tables. The general case, any finite list of atoms,
  is Theorem~\ref{thm:expcomplete}.
\end{theorem}

\begin{theorem}[Full deduction theorem for E]\label{thm:ddt}
  \uses{def:jops, thm:greedy}\lean{V.eimp_spec, V.ddt_E}\leanok
  The definable external implication
  $E(p,q)=\neg(p\wedge p)\vee(q\wedge q)$ satisfies
  $\Gamma,\varphi \vDash \psi \iff \Gamma \vDash E(\varphi,\psi)$ in
  both directions, over the whole language; the primitive arrow stays
  one-way. (This $E$ is a connective, not the judge's disposition E.)
\end{theorem}

\begin{theorem}[Blok--Pigozzi witnesses]\label{thm:bp}
  \uses{def:jops}\lean{V.dEq_spec, V.bp_refl, V.bp_mp, V.bp_cong_xnor, V.truth_equation, V.bp_iv}\leanok
  The same-value detector $\Delta$ and the truth equation
  $p\wedge p \approx \neg(p\wedge\neg p)$ satisfy conditions
  (i)--(iv) of the Blok--Pigozzi characterization on the matrix:
  ZTL is algebraizable (kinship: Bochvar algebras, Bonzio--Pra Baldi
  2024).
\end{theorem}

\begin{theorem}[Not self-extensional]\label{thm:selfext}
  \uses{def:jops}\lean{V.idem_interderivable, V.not_selfextensional, V.equiv_not_identical}\leanok
  $p \dashv\vDash p\wedge p$, yet $\neg(p\wedge p) \nvDash \neg p$:
  interderivability is not a congruence --- the same failure that
  builds the detectors.
\end{theorem}

\begin{theorem}[Structurality]\label{thm:structural}
  \uses{thm:greedy}\lean{V.evalF_subst, V.entails_structural}\leanok
  The substitution lemma: $\Gamma\vDash\varphi$ implies
  $\sigma\Gamma\vDash\sigma\varphi$; together with reflexivity,
  monotonicity and cut, $\vDash$ is a structural Tarskian consequence
  relation.
\end{theorem}

\begin{theorem}[Craig interpolation]\label{thm:interp}
  \uses{thm:expcomplete}\lean{V.Cpl.ztl_interpolation}\leanok
  If $\varphi \vDash \psi$, there is $\chi$ whose atoms all occur in both,
  with $\varphi\vDash\chi$ and $\chi\vDash\psi$: the J-DNF, over the shared
  atoms, of ``some value of $\varphi$'s private atoms makes $\varphi$ true''.
  With no shared atom it is a constant.
\end{theorem}

\chapter{Proof theory and quantifiers}

\begin{theorem}[Cut admissibility]\label{thm:cut}
  \uses{thm:closesiff}\lean{V.cut_admissible, V.weakening_admissible, V.identity_refutable}\leanok
  The tableau engine read bottom-up is a cut-free sequent calculus;
  cut on the covering pair $\{T\},\{F,Z\}$ is admissible (semantic cut
  elimination on top of the engine certificate), weakening admissible,
  identity derivable.
\end{theorem}

\begin{theorem}[Parameter tableaux, arbitrary domains]\label{thm:fo}
  \uses{thm:closesiff}\lean{ZParamEngine.entails_of_closed, ZParamEngine.entails_of_closed_core, ZParamEngine.run_fallen_bridge, ZParamHintikka.stuck_refutes}\leanok
  The $\gamma/\delta$ parameter rules carry the sign discipline to
  arbitrary domains. Kernel-checked: a closed run of the search without the
  classical step $\delta_2$ is a proof of entailment in every total model
  (with $\delta_2$, under its soundness as a stated hypothesis); without
  $\delta_2$ the search declines the fallen bridge
  $\neg\forall x P \vDash \exists x\neg P$; on closed inputs, a run that stops
  open yields a total countermodel read off the open branch, whose
  weakly signed atoms carry the mark. The soundness of $\delta_2$ needs
  $\neg\forall\to\exists\neg$; it is proved classically outside the
  zero-axiom corpus and enters it only as that hypothesis. Argued, not
  formalized: the infinite half of completeness (fair strategy, K\"onig's
  lemma, a model over an infinite domain). Undecidability of first-order ZTL (via the J-guard embedding of classical
  FO) is argued.
\end{theorem}

\begin{theorem}[Quantifier tableaux, finite domains]\label{thm:quant}
  \uses{thm:closesiff, def:jops}\lean{V.cover_allF_T, V.cover_allF_F, V.cover_exF_T, V.cover_exF_F, V.ui_mem, V.eg_mem, V.tab_ui, V.tab_drinker_needs_ground}\leanok
  Finite-domain quantifiers are strict folds in the certified language
  (a singleton domain collapses both to the $J_T$ guard); the $n$-ary
  signed rules are preimage-coverage theorems; UI/EG hold in membership
  form; eight battery verdicts are kernel evaluations of the certified
  engine.
\end{theorem}

\chapter{Grounding, general}

\begin{theorem}[Quarantine passport]\label{thm:passport}
  \uses{thm:kleene, thm:dynamics}\lean{V.passport_even2}\leanok
  Per-component diagnosis of ungroundedness: PARADOX (no classical
  models --- refusal permanent), UNDERDETERMINED (models exist ---
  liftable by stipulation), INPUT, DOWNSTREAM with culprits; the
  stipulation theorem measured totally; parity re-derived $62/62$.
  The even two-cycle's two classical models are kernel-checked.
\end{theorem}

\begin{theorem}[Knaster--Tarski for the lazy jump]\label{thm:kt}
  \uses{thm:kleene}\lean{V.evalK_mono, V.jumpL_mono, V.kt_fixed, V.kt_least, V.grounded_absolute}\leanok
  For every finite system: lazy evaluation is monotone over the whole
  language, the jump is monotone, the iteration from $\bot$ stabilizes
  within $n{+}1$ steps at the least fixed point, and the grounded part
  is identical in every fixed point --- quarantine is well-defined.
  Zero axioms (information-measure pigeonhole, no WF machinery).
\end{theorem}

\chapter{Expedition twins}

\begin{theorem}[Expedition twins]\label{thm:exped}
  \uses{thm:memz, def:jops}\lean{V.eqStream_never_T, V.eqStream_apart, V.diag_not_member, V.inj_cert_marked, V.mark_self_not_earned, V.atom_thresholds}\leanok
  Streams never earn identity while apartness is earned by a finite
  witness and persists; the diagonal earns strict non-membership; one
  marked pair collapses the injectivity certificate; under the separate
  reading of two sides (two acts of measurement) a nondegenerate mark does
  not equal itself; atom verdicts are the modal thresholds.
\end{theorem}

\begin{theorem}[Russell grounded, kernel-computed]\label{thm:russellkt}
  \uses{thm:kt}\lean{V.russell_grounded, V.russell_verdicts}\leanok
  With the constants $\top/\bot$ in the certified language, the lazy
  lfp of the nine-fact membership system computes to
  $[F,F,T,F,T,F,F,F,Z]$: eight facts grounded, quarantine exactly at
  $R\in R$ --- containment instead of explosion, by the kernel.
\end{theorem}

\begin{theorem}[Clean sets: C-extension]\label{thm:cleansets}
  \uses{thm:marked}\lean{V.memL_self, V.subL_refl_clean, V.seteq_refl_clean}\leanok
  On mark-free lists a verified element certifies its membership and
  $S \subseteq S$, $S = S$ hold --- the inheritance boundary is drawn
  at the mark, kernel-checked from both sides.
\end{theorem}
